
\documentclass[11pt]{article}
\usepackage[letterpaper,margin=0.72in]{geometry}
\usepackage[T1]{fontenc}
\usepackage[utf8]{inputenc}
\usepackage{lmodern}
\usepackage{array}
\usepackage{booktabs}
\usepackage{longtable}
\usepackage{tabularx}
\usepackage{enumitem}
\usepackage{xcolor}
\usepackage{hyperref}
\usepackage{fancyhdr}
\usepackage{lastpage}
\usepackage{pdflscape}
\usepackage{listings}
\hypersetup{colorlinks=true,linkcolor=black,urlcolor=blue,citecolor=black}
\definecolor{HeaderBlue}{HTML}{1F4E79}
\definecolor{LightBlue}{HTML}{EAF3F8}
\definecolor{LightGray}{HTML}{F5F7FA}
\definecolor{RuleGray}{HTML}{B7C3D0}
\definecolor{RiskAmber}{HTML}{FFF4D6}
\newcolumntype{Y}{>{\raggedright\arraybackslash}X}
\newcolumntype{L}[1]{>{\raggedright\arraybackslash}p{#1}}
\setlength{\parindent}{0pt}
\setlength{\parskip}{0.45em}
\setlist[itemize]{leftmargin=1.2em,itemsep=0.18em,topsep=0.2em}
\setlist[enumerate]{leftmargin=1.4em,itemsep=0.18em,topsep=0.2em}
\pagestyle{fancy}
\setlength{\headheight}{14pt}
\fancyhf{}
\lhead{Gauntlt DevSecOps Security Testing}
\rhead{\thepage\ of \pageref{LastPage}}
\renewcommand{\headrulewidth}{0.4pt}
\sloppy
\lstset{basicstyle=\ttfamily\small,breaklines=true,columns=fullflexible,keepspaces=true,frame=single,framerule=0.3pt,rulecolor=\color{RuleGray},backgroundcolor=\color{LightGray}}
\newcommand{\TitleBlock}[2]{%
  \begin{center}
  {\LARGE\bfseries\color{HeaderBlue} #1}\\[0.35em]
  {\large #2}\\[0.4em]
  {\small Prepared for DevSecOps implementation planning | Updated May 6, 2026}
  \end{center}
  \vspace{0.4em}\hrule\vspace{0.8em}
}
\newsavebox{\gauntltbox}
\newenvironment{Callout}[1]{%
  \par\vspace{0.3em}\noindent\begin{lrbox}{\gauntltbox}\begin{minipage}{0.96\linewidth}\textbf{#1}\par\smallskip}
  {\end{minipage}\end{lrbox}\fcolorbox{RuleGray}{LightBlue}{\usebox{\gauntltbox}}\par\vspace{0.4em}}
\newenvironment{RiskBox}[1]{%
  \par\vspace{0.3em}\noindent\begin{lrbox}{\gauntltbox}\begin{minipage}{0.96\linewidth}\textbf{#1}\par\smallskip}
  {\end{minipage}\end{lrbox}\fcolorbox{RuleGray}{RiskAmber}{\usebox{\gauntltbox}}\par\vspace{0.4em}}
\newcommand{\checkbox}{\raisebox{0.2ex}{\fbox{\rule{0pt}{0.9ex}\rule{0.9ex}{0pt}}}}
\begin{document}

\TitleBlock{Gauntlt Security Testing Handout}{Security-as-code attacks for CI/CD pipelines}

\begin{Callout}{Purpose}
This handout explains how to use Gauntlt as a lightweight security testing harness for CI/CD. It focuses on authorized, repeatable DAST-style checks that wrap common command-line tools and express expected security outcomes as readable attack files.
\end{Callout}

\section{What Gauntlt Is}
Gauntlt is an open-source, Ruby-based security testing framework built around the idea of being intentionally rigorous with deployed code before it reaches production. It uses Cucumber/Gherkin-style text files with the \texttt{.attack} extension. Each attack file describes security scenarios in language that can be reviewed by developers, security engineers, and operations staff.

Gauntlt is best understood as a \textbf{security test orchestrator}, not a replacement for a scanner. It wraps existing tools such as \texttt{nmap}, \texttt{curl}, \texttt{dirb}, \texttt{sqlmap}, \texttt{SSLyze}, Arachni, Garmr, and custom command-line scripts.

\section{Recommended Use Cases}
\begin{tabularx}{\linewidth}{@{}L{0.26\linewidth}Y@{}}
\toprule
\textbf{Use Case} & \textbf{How Gauntlt Helps} \\
\midrule
Pipeline smoke testing & Runs focused security checks after deployment to a test environment. \\
Port exposure validation & Uses \texttt{nmap} or equivalent tooling to confirm only expected ports are reachable. \\
TLS posture checks & Uses tools such as \texttt{SSLyze} to check protocol, certificate, and cipher configuration. \\
HTTP control validation & Uses \texttt{curl} to check headers, redirects, authentication behavior, and error handling. \\
Regression checks & Encodes known security requirements as pass/fail tests that run every build. \\
Security collaboration & Makes security tests readable through Gherkin syntax and version-controlled attack files. \\
\bottomrule
\end{tabularx}

\section{Where Gauntlt Fits in a DevSecOps Pipeline}
\begin{enumerate}
  \item Build and package the application.
  \item Deploy to an authorized ephemeral, staging, or security test environment.
  \item Run Gauntlt attack files against that environment.
  \item Store the output as CI evidence.
  \item Fail the pipeline when defined security expectations are not met.
  \item Route failures to the owning team with the attack file, command output, and remediation note.
\end{enumerate}

\section{Installation Pattern}
\begin{lstlisting}[language=bash]
# Install Ruby and Bundler using the approved platform method first.
gem install gauntlt

# Confirm the CLI is available.
gauntlt --help

# Run a single attack file.
gauntlt attacks/tls_headers.attack

# Run a directory of attack files.
gauntlt attacks/*.attack
\end{lstlisting}

\section{Example Attack File}
The following example is intentionally safe and focused on verifying an HTTP security header. Adapt targets, paths, and assertions to your approved test environment.

\begin{lstlisting}
Feature: Verify baseline web security headers

  Scenario: Application returns X-Content-Type-Options
    Given "curl" is installed
    When I launch a "curl" attack with:
      """
      curl -s -I https://example.internal/
      """
    Then the output should contain:
      """
      X-Content-Type-Options: nosniff
      """
\end{lstlisting}

\section{Operating Rules}
\begin{itemize}
  \item Run attacks only against systems where the team has written authorization.
  \item Keep attack files narrow, deterministic, and mapped to explicit security requirements.
  \item Use staging or ephemeral environments before considering any production-facing checks.
  \item Treat scanner output as evidence, but require human-readable failure messages for remediation.
  \item Pin tool versions where possible to avoid noisy pipeline drift.
  \item Start with non-destructive checks such as headers, TLS posture, expected ports, and authentication redirects.
\end{itemize}

\section{Legacy Tooling Note}
Gauntlt remains useful as a security-as-code pattern and as a simple wrapper for command-line security tools. For modern enterprise adoption, evaluate maintenance posture, Ruby compatibility, adapter reliability, and whether newer tools such as OWASP ZAP, Nuclei, Semgrep, Trivy, or custom policy-as-code checks better fit the current control objective.


\section*{References}
\begin{itemize}
  \item Gauntlt official site: \url{http://gauntlt.org/}
  \item Gauntlt GitHub repository: \url{https://github.com/gauntlt/gauntlt}
  \item Gauntlt Starter Kit: \url{https://github.com/gauntlt/gauntlt-starter-kit}
  \item Gauntlt Demo: \url{https://github.com/gauntlt/gauntlt-demo}
\end{itemize}

\end{document}



